\label{chap:83-06}

\section{Projective measure and evolutionary conservation of the unit}
\label{p12-83-c6-s1:chap:medida-proyectiva}

Cylinder topology provides a notion of proximity; a theory of paths also requires weights. The evolutionary conservation of the unit here takes an exact probabilistic form: every cylinder distributes its weight among its extensions, and the total sum remains equal to one. The unit is not fixed in a cell; it is transported through refinement.

\subsection{Consistent weights on finite levels}

Let \(\mu_n\) be a probability on the finite set
\(\mathcal S_n\). The family is \emph{projectively consistent} when
\begin{equation}
 \mu_n(a)
 =\sum_{\substack{b\in\mathcal S_{n+1}\\
                   \rho_{n+1,n}(b)=a}}
   \mu_{n+1}(b)
 \qquad(a\in\mathcal S_n).
 \label{p12-83-c6-s1:eq:consistencia-pesos}
\end{equation}
This equality expresses the local preservation of the unit: the weight of a
truncated history is the sum of the weights of all its children.

\begin{theorem}[Genealogical measure on cylinders]
\label{p12-83-c6-s1:thm:medida-genealogica}
Every family of finite probabilities that satisfies
\eqref{p12-83-c6-s1:eq:consistencia-pesos} determines a unique Borel probability
\(\mu_\infty\) on \(\Omega_\infty\) such that
\begin{equation}
 \mu_\infty([a]_n)=\mu_n(a).
 \label{p12-83-c6-s1:eq:medida-cilindro}
\end{equation}
\end{theorem}

\begin{proof}
\[
 \mathscr A
 :=\{\text{finite unions of cylinders}\}
\]
is an algebra: the intersection of two cylinders is either empty or a cylinder of the
deeper level. Every element of \(\mathscr A\) can be refined into a
disjoint union of cylinders of a common level. If
\(A=\bigsqcup_{a\in I}[a]_n\), we set
\[
 \mu_0(A)=\sum_{a\in I}\mu_n(a).
\]
Iterating \eqref{p12-83-c6-s1:eq:consistencia-pesos} proves that the value does not depend on the
chosen common level.

If a cylinder \([a]_n\) is empty, finite branching and Koenig's lemma
\cite{Koenig1936}
imply that there is a depth with no descendants of \(a\). Iterated
consistency then gives \(\mu_n(a)=0\). Hence \(\mu_0\) is well
defined even for representations containing empty cylinders and is
finitely additive.

The cylinders are both open and closed. If a decreasing sequence of
elements of \(\mathscr A\) had empty intersection, compactness
would imply that one of them was already empty; hence \(\mu_0\) is continuous at
the empty set and constitutes a countably additive premeasure. The Caratheodory extension theorem ---equivalently, the Kolmogorov extension for this finite
projective system--- extends it to the \(\sigma\)-algebra generated by the
cylinders \cite{Caratheodory1914,Kolmogorov1933}. This is the Borel algebra of
\(\Omega_\infty\), and finiteness of the
measure ensures uniqueness.
\end{proof}

The measure may follow from local transitions. If
\(P_n(b\mid a)\ge0\) satisfies
\[
 \sum_{\rho_{n+1,n}(b)=a}P_n(b\mid a)=1,
\]
then
\(\mu_{n+1}(b)=\mu_n(a)P_n(b\mid a)\) is consistent. Uniformity is a
possible choice, not an obligation. The weights may depend on the sheet,
carry, orientation, a nonnegative action cost, or a record, provided
they preserve the sum. For example,
\[
 P_n(b\mid a)
 =\frac{e^{-\beta a_n(b\mid a)}}
        {\sum_{b'\mapsto a}e^{-\beta a_n(b'\mid a)}},
 \qquad a_n\geq0,
\]
defines normalized positive weights. The orthogonal action phase belongs to the
kernel of \cref{chap:integral-genealogica}; it is not inserted into these
positive probabilities.

\subsection{Decomposition of the genealogical measure on fibers of equal Archimedean value}

The Archimedean evaluation transports the measure to \(K\):
\begin{equation}
 \nu:=\Val_*\mu_\infty,
 \qquad
 \nu(B)=\mu_\infty(\Val^{-1}(B)).
 \label{p12-83-c6-s1:eq:pushforward-arquimediano}
\end{equation}
The visible probability of a region sums the weight of all the genealogies that
publish it. When two sections possess the same value, their contributions are
aggregated in \(\nu\), although they remain separate in \(\mu_\infty\).

\begin{proposition}[Value Decomposition]
If \(K\) is a standard Borel space, there exists a family of conditional
measures \(\{\mu_y\}_{y\in K}\), defined for \(\nu\)-almost every \(y\),
such that \(\mu_y\) is supported on \(\mathfrak G_y=\Val^{-1}(y)\) and
\[
 \int_{\Omega_\infty}F\dd\mu_\infty
 =\int_K\left(\int_{\mathfrak G_y}F\dd\mu_y\right)\dd\nu(y)
\]
for every integrable function \(F\).
\end{proposition}

This formula separates value and genealogy at the level of the measure. Physics that
depends only on the value first integrates the entire fibre; a memory observable
can distinguish different distributions \(\mu_y\) over the same
real coordinate.

\subsection{Conditioning as pruning of futures}

Let \(A\subseteq\Omega_\infty\) be an observable event with
\(\mu_\infty(A)>0\). The conditional measure is
\begin{equation}
 \mu_\infty(B\mid A)
 =\frac{\mu_\infty(B\cap A)}{\mu_\infty(A)}.
 \label{p12-83-c6-s1:eq:condicionamiento-genealogico}
\end{equation}
When \(A\) is the reading of a register, its intersection with each cylinder
eliminates the incompatible extensions. Normalization does not create new branches:
it redistributes the unit among those that have survived.

\begin{proposition}[Compatible update]
If \(A\) is a union of cylinders at level \(m\), then for every \(n\ge m\)
the conditional distribution on \(\mathcal S_n\) is obtained by removing the
states whose truncations do not belong to \(A\) and renormalizing the remaining
weights. The conditional distributions continue to satisfy
\eqref{p12-83-c6-s1:eq:consistencia-pesos}.
\end{proposition}

\begin{proof}
The identity is Bayes' rule on finite sets. The sum of the
weights of the conditioned children coincides with the conditioned weight of the parent,
because intersection and disjoint union commute.
\end{proof}

\subsection{Combinatorial conservation and dimensional completion}

The same law appears in the binomial completion:
\[
 1=\sum_{r=0}^{d}\binom dr
 \left(\frac9{10}\right)^{d-r}
 \left(\frac1{10}\right)^r.
\]
In three dimensions,
\[
 1000=729+243+27+1,
\]
so that the interior, new faces, extreme edges, and final vertex form
a partition of decimal unity. The passage
\(729+270=999\to729+271=1000\) distinguishes the state prior to closure of the
completed unity. In the path system, the analogous identity is
\[
 \mu([a]_n)=\sum_{b\mapsto a}\mu([b]_{n+1}).
\]
The unit evolves as it is distributed among the strata and is recovered by summing
the complete genealogy.

\begin{figure}[htbp]
\centering
\begin{tikzpicture}[>=Latex,node distance=12mm and 17mm,
  box/.style={draw=hmtblue,rounded corners,fill=hmtlightblue,align=center,
              minimum width=28mm,minimum height=8mm},
  leaf/.style={draw=hmtgreen,rounded corners,fill=hmtlightgreen,align=center,
               minimum width=22mm,minimum height=7mm}]
\node[box] (a) {$[a]_n$\\weight $\mu_n(a)$};
\node[leaf,below left=of a] (b1) {$[b_1]_{n+1}$\\$\mu_{n+1}(b_1)$};
\node[leaf,below=of a] (b2) {$[b_2]_{n+1}$\\$\mu_{n+1}(b_2)$};
\node[leaf,below right=of a] (b3) {$[b_3]_{n+1}$\\$\mu_{n+1}(b_3)$};
\draw[->] (a)--(b1); \draw[->] (a)--(b2); \draw[->] (a)--(b3);
\node[below=20mm of b2,align=center,text width=.78\textwidth] (eq)
 {$\mu_n(a)=\mu_{n+1}(b_1)+\mu_{n+1}(b_2)+\mu_{n+1}(b_3)$:\\
the unit is preserved when the genealogy is refined.};
\end{tikzpicture}
\caption{Projective weight conservation on cylinders.}
\label{p12-83-c6-s1:fig:conservacion-peso-cilindros}
\end{figure}

\begin{statusbox}
The projective measure is constructed on the HMT inverse system.
The extension theorem is classical mathematics; the identification of its
cylinders with TPK states is exact once the transition weights have been fixed.
The selection of a specific physical measure belongs to the dynamical model and does not
follow solely from the census of branches.
\end{statusbox}

The projective measure is characterized by positivity, conservation of
weight between parents and children, and conditioning within the observed event. A
complete experimental realization adds preparation and the sampling
protocol; with those data, the same measure produces reading probabilities
comparable with frequencies.
